% !TeX root = ../papers/nc-p3.tex
\section{四中心切断を持つ全有限窓の同型不変量}
\label{sec:central-cut-kummer-classification}
\begin{theorem}\label{thm:central-cut-kummer-classification}
前節の二つの中心型の全ての窓に対し $U=X_{\rm sm}$ とし、
$\beta\in H^2_{\rm et}(U,\mu_s)$ を、零番の次数 $s$ のAzumaya隅の
$\mathrm{PGL}_s$ 束を $\mathrm{SL}_s$ へ持ち上げる障害とする（符号は前節までの行列式規約）。
二窓が添字保持同型であるための必要十分条件は、偏極中心間に $\beta$ を運ぶ同型が存在することである。
同じ $(X,H)$ の上の印付き同型類は有限である。
$U$ 上の直線束 $L$ で $4L=sH$ となり、一移動で $\beta$ は $[L]\bmod s$ だけ増える。
実際 $s=4$ なら $L=H$、$s=2$ なら $2L=H$ である。反対環を同じ始点から取れば
印は $-\beta$ となる。特に $s=2$ の窓はすべて周期が2の約数である。
\end{theorem}
\begin{proof}[一般切断の因子と分裂束]
$K=k(|H|)$、$S/K$ を一般超平面、$S^\circ=S\cap U$ とする。
制限は $\operatorname{Pic}(U)\simeq\operatorname{Pic}(S^\circ)$ である。
実際 $U$ 上の普遍接続 $J$ は射影空間束で、そのPicard群は $\operatorname{Pic}(U)$ と
パラメータ射影空間からの超平面束の直和である。
一般の二超平面の共通切断は滑らかな整曲線で $U$ に含まれ、
そこで二切断は横断的である（豊富な基点無し系のBertiniによる）。
この二切断の鉛筆では全てのメンバーが整である。
各因子成分は豊富性により基底曲線と曲線で交わる、すなわち基底曲線を含むが、
その一般点でメンバーは滑らかだからである。Cartier切断はCohen--Macaulayなので埋込み成分もない。
よって幾何学的非整切断の閉集合は余次元2以上である。
普遍族は平坦だから垂直素因子は底の因子を支配し、その一般繊維は幾何学的に整で
引き戻しの重複度は1である。正則な $J$ に一般繊維の因子を閉包として延長すれば
Picard群の制限核はパラメータ空間からの因子で生成され、主張を得る。

$\mathcal S$ を $S$ の標準滑らかなスタックとする。滑らかな場合は $S$ 自身で、
幾何的な各 $A_1$ チャートは $[\mathbb A^2/\{\pm1\}]$ のエタール局所型である。
商特異点の標準スタックは滑らかな部分の恒等写像を延長する標準構成なので $K$ 上で取れる。
$\mu_s$ 係数の第二コホモロジーは正則な曲面で余次元2を除いても変わらない
（絶対純性をエタールチャートで適用）。従って $S^\circ$ 上の分裂の $\mu_s$-gerbeを延ばせる。
ねじれ束の反射的延長は滑らかな曲面上では局所自由である。
普遍接続の関数体は $k(X)$ の純超越拡大なので、$K(S)$ 上の指数は依然 $s$ である。
四束 $E_i$ は従って $K$ 上で最小の正階数を持つ。$E_{i+4}=E_i(h)$、$h=H|_S$ とする。
二束間のHomの粗曲面への直像は元の隅の制限である。両方反射的で、
前節の局所型と極大Cohen--Macaulay性により穴あき部分から一意に延びるからである。
有限安定群の不変部分は完全関手なのでEuler数も一致する。

非自明な安定群の上では、局所ねじれを自明化した後の各束は自明・符号表現を二つずつ含む。
よって非恒等元のtraceは零であり、自明化を取り替えてもこれは変わらない。
任意の二つのこの型の束 $E,F$ の $\mathcal Hom(E,F)$ についても同じである。
従って滑らかな商スタックのRiemann--Roch（To\"enの定理、複素数上では軌道体の指数公式）で、
非自明な慣性成分の補正は全て零となる。
$g=c_1(\det F/\det E)$ とすると、恒等成分の通常のChern指標による計算から
\[
 \chi(\mathcal Hom(E,F))
 =\tfrac12\{\chi(\mathcal End E)+\chi(\mathcal End F)+g^2\}.
\]
実際 $D=\mathcal Hom(E,F)$ について
$\mathcal End(D)=\mathcal End E\otimes\mathcal End F$、$c_1(D)=s g$ なので、
第二指標の比較は $2\operatorname{ch}_2D=
\operatorname{ch}_2(\mathcal End E)+\operatorname{ch}_2(\mathcal End F)+g^2$ になる。
標準束は自明で、階数に比例するTodd項は両辺で一致する。

窓内で $F=E_j$, $E=E_i$, $a=j-i$ とする。
切断Hilbert多項式 $32m^2+16am+2a^2+2$ から
$g^2=4a^2$, $sgh=16a$, $h^2=64/s^2$ を得る。
行列式の比は各安定群で自明である（同じ重複度の表現なので行列式は一致する）。
従ってこれは粗曲面上の直線束になる。その最小解消はK3であり、Hodge指数定理の等号から
$g=(s/4)a h$ が数値的に従う。K3上で数値同値と線型同値が一致し、粗曲面からの
引き戻しも単射なので、四倍した式は直線束の等式である。
幾何的基礎拡大から $K$ へも、固有幾何学的整曲面のPicard群の単射で降りる。
特に全ての連続行列式比は同じ $L$ になり、$s=4$ で $L=H$、$s=2$ で $2L=H$。
$S^\circ$ から $U$ へは冒頭のPicard同型を使う。
\end{proof}

\begin{proof}[二超平面における同型の降下]
同じ中心上で $\beta$ が等しければBrauer類が等しく、二窓の束を同じgerbe上に置ける。
Kummer列によって零番の行列式の比は $s$ 乗である。
一つの通常の直線束で全束を同時に捩れば両列の零番行列式は一致する。
連続差 $L$ も両者で一致する。$s=2$ の場合も、二つの半分は
一般切断上のK3で一致し、上のPicard同型で $U$ に戻る。
対応する各二束で $g=0$、自己HomのEuler数は2なので上式から相互HomのEuler数も2である。
$\mathcal S$ の標準束は自明である。Serre双対により一方向に非零射が存在し、
最小指数により一般階数は $s$ である。行列式の次数が同じなので余核に因子台はなく、
同階数の束の行列式は全点で可逆となる。よって $K$ 上で各対は同型である。

一般の二超平面を独立に選び、その共通曲線 $C$ を対の定義体上で取る。
$H^0(C,\mathcal End E_i)$ はその定義体だけである。
実際経路長零の対角隅の次数零切断は定数で、二超平面のKoszul列の核・中間項からは
新しい切断は出ない。隅が深さ4で全捩りの第一・第二コホモロジーを持たず、
負の経路長の切断がないからである。従って二つの一般切断の同型は共通曲線上では倍率だけ異なる。
その張るHom空間の一次元部分は $J\times_U J$ の一般点で二つの引き戻しが一致する。
よって $U$ の一般点で定まる。例えば $J\to U$ は射影空間束なので
二組の独立な不定元で値が等しい射影座標比は底の関数体に入る。
この有理直線を $U$ の二束のHom束内で飽和し二重双対を取って直線束 $P_i$ とする。
評価 $E_i\otimes P_i\to E'_i$ は反射性で全 $U$ へ延びる。
各高さ1では飽和生成元の成分に単元がある。その因子を一般超平面で切ると、
同じ射影線は真の束同型で生成されるので行列式もそこでは単元である。
従って行列式に零因子はなく評価は全 $U$ で同型である。
$S^\circ$ 上の直線 $P_i$ は選んだ束同型で自明である（余次元1で一致すれば足りる）。
Picard同型から $P_i$ は自明である。こうして四束は最初の共通捩りだけを除き
$U$ 上で実際に一致する。

各隅は $X$ 上反射的で、$U$ 上の束Homから一意に戻る。
中心次数の全切断も深さ4により元の隅である。よって得た束同型から中心同定を保って
全ての経路と積を復元する。逆に窓の同型は標準折り畳みにより中心と零番の
Azumaya束を対応させるため $\beta$ を運ぶ。
固定中心上の候補は有限群 $H^2(U,\mu_s)$ の部分に入る。
有限性は複素数上の有限型多様体の有限係数コホモロジーの有限性と幾何学的な比較による。
移動の公式は行列式比の計算、反対の公式は双対束を取ることから従う。
\end{proof}
